\chapter{Conclusiones y propuestas}
\label{cap:conclusiones}

Este último capítulo recoge las conclusiones extraídas del trabajo,
valorando el grado de cumplimiento de los objetivos planteados en el
apartado~\ref{sec:objetivos}, y propone líneas de ampliación futuras del laboratorio.

\section{Conclusiones}

A continuación se retoma cada uno de los objetivos específicos planteados
en el apartado~\ref{sec:objetivos}, valorando su grado de cumplimiento a la vista del
trabajo realizado:

\begin{itemize}
\item \textbf{Estudiar una metodología de referencia existente en el ámbito
  de la seguridad de aplicaciones web (OWASP WSTG).} Cumplido en el
  capítulo~\ref{cap:marco-teorico}, donde se analiza la estructura, fases y
  catálogo de pruebas de la WSTG y se justifica su elección frente a otros
  marcos como PTES, OSSTMM o NIST~SP~800-115.
\item \textbf{Diseñar un entorno web vulnerable controlado.} Cumplido con el
  diseño e implementación de \breachlab, descrito en el
  capítulo~\ref{cap:diseno}, incluido su despliegue reproducible mediante
  Docker Compose.
\item \textbf{Derivar, a partir de la metodología estudiada, una
  metodología propia orientada a la agilidad.} Cumplido en el
  capítulo~\ref{cap:metodologia}, acotando el alcance de la WSTG a un
  subconjunto de seis pruebas alineadas con el OWASP Top~10 y consolidando
  sus fases en un ciclo único de seis pasos, pensado para equipos con
  recursos limitados.
\item \textbf{Identificar vulnerabilidades siguiendo dicha metodología
  propia.} Cumplido: cada apartado del capítulo~\ref{cap:explotacion}
  localiza la vulnerabilidad en una línea de código concreta antes de
  explotarla, siguiendo el proceso del capítulo~\ref{cap:metodologia}.
\item \textbf{Explotar las vulnerabilidades de forma controlada.} Cumplido
  en un entorno aislado, sin exposición a Internet, con usuarios y datos
  ficticios, y con evidencia registrada de cada explotación en
  \texttt{backend/logs/attacks.log}.
\item \textbf{Analizar el impacto técnico y funcional de cada
  vulnerabilidad.} Cumplido mediante la valoración CVSS v3.1 de cada
  vulnerabilidad (capítulo~\ref{cap:explotacion}) y, de forma adicional a lo
  previsto inicialmente, mediante el caso de estudio encadenado del
  capítulo~\ref{cap:caso-estudio}, que muestra cómo ese impacto se
  amplifica cuando varias vulnerabilidades se combinan.
\item \textbf{Comparar los resultados y el esfuerzo de la metodología
  propia frente a la metodología de referencia.} Cumplido en el
  capítulo~\ref{cap:comparativa}, que contrasta la cobertura y el esfuerzo
  estimado de la metodología propia frente a la OWASP WSTG completa, y
  complementa la comparación situando \breachlab frente a DVWA y OWASP
  Juice Shop. Esta comparación de esfuerzo es cualitativa y no una medición
  cronometrada (apartado~\ref{sec:comparativa-metodologica}), lo que se
  recoge como limitación del trabajo y como línea futura.
\end{itemize}

\section{Recomendaciones de mitigación}

A modo de síntesis práctica orientada a quien despliegue una aplicación
similar en producción, la tabla~\ref{tab:mitigaciones} recoge, para cada
vulnerabilidad tratada, la recomendación principal aplicada en la rama
segura de \breachlab y la \emph{cheat sheet} de la OWASP Foundation
correspondiente.

\begin{table}[htbp]
  \centering
  \small
  \begin{tabular}{lp{6.2cm}l}
    \toprule
    \tabhead{Vulnerabilidad} & \tabhead{Recomendación principal} & \tabhead{Referencia} \\
    \midrule
    Inyección SQL &
      Usar siempre consultas parametrizadas, nunca concatenar entrada de usuario &
      \citelink{OWASPCS-SQLi} \\
    XSS (DOM) &
      Insertar datos no confiables como texto (\texttt{textContent}), nunca como HTML, y aplicar CSP &
      \citelink{OWASPCS-XSS} \\
    CSRF &
      Token impredecible ligado a la sesión, verificado en el servidor antes de cambios de estado &
      \citelink{OWASPCS-CSRF} \\
    Fallos de autenticación &
      Hash de contraseñas, bloqueo tras varios intentos fallidos y mensajes de error genéricos &
      \citelink{OWASPCS-AUTH} \\
    IDOR &
      Verificar la propiedad del recurso (o el rol) en el servidor en cada acceso, nunca confiar en el identificador recibido &
      \citelink{OWASPCS-AUTHZ} \\
    Configuración insegura &
      Cabeceras de hardening (CSP, \texttt{X-Frame-Options}, \texttt{X-Content-Type-Options}) activas por defecto &
      \citelink{OWASPCS-HEADERS} \\
    \bottomrule
  \end{tabular}
  \caption{Recomendaciones de mitigación aplicadas en \breachlab, con la
    \emph{cheat sheet} de la OWASP Foundation correspondiente a cada una.}
  \label{tab:mitigaciones}
\end{table}

\section{Trabajo futuro y posibles ampliaciones}
\label{sec:trabajo-futuro}

Para cerrar este último capítulo, se recogen algunas
propuestas que podrían añadirse al trabajo actual en el futuro:

\begin{itemize}
\item Medir de forma cuantitativa (tiempo de aplicación, número de hallazgos
  por hora) la diferencia de esfuerzo entre la metodología propia y la
  OWASP WSTG completa, aplicando ambas sobre \breachlab y sobre algún
  entorno externo adicional, para sustituir la estimación cualitativa del
  apartado~\ref{sec:comparativa-metodologica} por datos medidos.
\item Ampliar la cobertura de vulnerabilidades de \breachlab a otras
  categorías del OWASP Top~10 no tratadas en este trabajo, como
  deserialización insegura, SSRF o componentes con vulnerabilidades
  conocidas.
\item Añadir variantes adicionales de las vulnerabilidades ya existentes:
  una inyección SQL ciega basada en tiempo en un buscador de notas, y un XSS
  reflejado además del XSS basado en DOM ya implementado.
\item Sustituir los identificadores secuenciales de las notas por
  identificadores UUID no predecibles, para discutir en qué medida esto
  mitiga el IDOR frente a una comprobación real de propiedad.
\item Automatizar la comparativa del capítulo~\ref{cap:comparativa} mediante
  un escaneo DAST (por ejemplo, OWASP ZAP \emph{baseline}) contra el modo
  vulnerable y el modo \texttt{?safe=1}, incorporando sus resultados como
  tabla objetiva.
\item Incorporar autenticación multifactor y políticas de expiración y
  rotación de sesión como contramedida adicional frente a la cadena de
  ataque descrita en el capítulo~\ref{cap:caso-estudio}.
\item Persistir el contador de intentos fallidos de \verb|login()|
  (actualmente en memoria del proceso y por nombre de usuario) en un
  almacén compartido, de forma que sobreviva a un reinicio del backend y
  permita además detectar fuerza bruta distribuida por dirección IP.
\item Registrar también en \texttt{attacks.log} el propio bloqueo temporal
  (HTTP 429) del apartado~\ref{sec:auth}, que actualmente solo queda
  reflejado en el log de acceso HTTP de Werkzeug y no mediante
  \verb|log_event()|, a diferencia del resto de eventos de seguridad de la
  aplicación.
\end{itemize}
